% !TEX root = ../main.tex
\begin{ouverture}[boucles]{14}{Automatismes et synthèses}{Partie 5 · Méthodes transversales}
{Calcul\enspace·\enspace Pourcentages\enspace·\enspace Lectures graphiques\enspace·\enspace Problèmes transversaux}
\objectif{Entretenir les calculs et les lectures de représentations.}
\objectif{Choisir rapidement une méthode adaptée.}
\objectif{Combiner plusieurs chapitres dans une même résolution.}
\objectif{Rédiger et contrôler une réponse complète.}
\prerequis{Les exercices 1 à 8 réactivent les acquis antérieurs.}
\prerequis{Les synthèses finales mobilisent les chapitres du livre.}
\end{ouverture}
\essentiel
\essSection{1}{Des rituels à répartir sur l'année}
\begin{savoirfaire}
Les exercices courts de ce chapitre sont une banque de réactivation.
Travailler régulièrement quelques questions de thèmes différents, avant
de passer aux problèmes. Les automatismes se construisent tout au long de
l'année ; ils ne constituent pas une leçon à apprendre en une seule séance.
\end{savoirfaire}
\begin{propriete}[Évolutions en pourcentage]
Une évolution de taux $t$ multiplie par $1+t$. Des évolutions successives
multiplient leurs coefficients : le taux global est le produit moins 1.
Pour revenir à la valeur initiale après un coefficient $c>0$, utiliser $1/c$.
Un taux réciproque n'est donc pas en général l'opposé du taux initial.
\end{propriete}
\begin{exemple}
Une hausse de 25\,\% suivie d'une baisse de 20\,\% donne
$1{,}25\times0{,}8=1$ : retour à la valeur initiale.
Après une baisse de 20\,\%, il faut une hausse de $1/0{,}8-1=25\,\%$ pour revenir.
\end{exemple}
\essSection{2}{Calculer efficacement}
\begin{propriete}
Identités : $(a+b)^2=a^2+2ab+b^2$, $(a-b)^2=a^2-2ab+b^2$,
$a^2-b^2=(a-b)(a+b)$. Un produit est nul si au moins un facteur est nul.
Avant de développer, identifier l'objectif : racines, signe, dérivée ou valeur numérique.
\end{propriete}
\begin{attention}
Une fraction ne se simplifie que par des facteurs communs non nuls.
L'ensemble des solutions doit respecter les valeurs interdites et le domaine du modèle.
\end{attention}

\clearpage
\essSection{3}{Lire et produire des représentations}
\begin{propriete}
La droite passant par $A(x_A;y_A)$ et $B(x_B;y_B)$, avec $x_A\neq x_B$,
a pour pente $m=(y_B-y_A)/(x_B-x_A)$. Son équation est $y=mx+p$,
avec $p=y_A-mx_A$. Dans un diagramme circulaire, une fréquence $f$ correspond
à un angle de $360f$ degrés.
\end{propriete}
\begin{savoirfaire}
Sur une courbe, résoudre $f(x)=k$ revient à lire les intersections avec la
droite horizontale $y=k$ ; pour $f(x)\geq k$, lire les portions au-dessus.
Contrôler les unités, l'origine et l'échelle de chaque axe. Dans un histogramme,
les fréquences sont portées par les aires ; dans une boîte, les repères sont les quartiles.
\end{savoirfaire}
\essSection{4}{Rédiger un problème de synthèse}
\begin{savoirfaire}
Définir les objets et le domaine. Écrire la méthode avant le calcul.
Présenter les étapes et vérifier les bornes. Donner une réponse dans les unités
de la question et distinguer valeur exacte, approximation et conjecture.
Une étude de signe, un tableau de variations ou un arbre peut soutenir l'explication.
\end{savoirfaire}
\begin{attention}
Les quatre derniers exercices sont des sujets originaux de synthèse, pas
des reproductions de sujets officiels. Ils mobilisent plusieurs chapitres.
Chercher sans le corrigé, puis reprendre les notions signalées par les erreurs.
\end{attention}

\exercices

\rubrique{Découverte}
\begin{blocLivre}{ex-14-1}
\begin{exercice}[1]{Évolutions et retour au départ}
Un prix de 80 euros augmente de 25\,\%, puis diminue de 20\,\%.\begin{questions}
\item Calculer le prix après chaque évolution et le taux global.
\item Quel taux de hausse compense exactement une baisse de 20\,\% ?
\end{questions}
\end{exercice}
\end{blocLivre}
\begin{blocLivre}{ex-14-2}
\begin{exercice}[1]{Retrouver une valeur et un taux}
\begin{questions}
\item Un article soldé de 30\,\% coûte 84 euros. Quel était son prix initial ?
\item Une quantité passe de 45 à 54. Calculer son taux d’évolution.
\item Calculer le taux global de deux hausses successives de 10\,\% puis 20\,\%.
\end{questions}
\end{exercice}
\end{blocLivre}
\begin{blocLivre}{ex-14-3}
\begin{exercice}[1]{Développer, factoriser, résoudre}
\begin{questions}
\item Développer $(2x-3)^2$.
\item Factoriser $9x^2-25$.
\item Résoudre $(x-2)(x+1)=0$.
\end{questions}
\end{exercice}
\end{blocLivre}
\begin{blocLivre}{ex-14-4}
\begin{exercice}[1]{Signe et intervalles}
\begin{questions}
\item Résoudre $3-2x\geq0$.
\item Résoudre $(x+2)(x-3)<0$.
\item Donner un exemple de solution commune aux deux inéquations.
\end{questions}
\end{exercice}
\end{blocLivre}
\begin{blocLivre}{ex-14-5}
\begin{exercice}[1]{Équations de droites}
\begin{questions}
\item Déterminer la droite passant par $A(1;3)$ et $B(3;7)$.
\item Pour $d:y=-x/2+4$, donner pente, ordonnée à l’origine et intersection avec l’axe des abscisses.
\end{questions}
\end{exercice}
\end{blocLivre}
\clearpage
\rubrique{Entraînement}
\begin{blocLivre}{ex-14-6}
\begin{exercice}[2]{Lire une courbe par morceaux}
La fonction $f$ est représentée par les segments suivants ; les sommets ont pour coordonnées exactes $(-2;1)$, $(0;-1)$, $(2;3)$, $(4;1)$.\begin{center}
\begin{tikzpicture}\begin{axis}[otLivreGraphe,width=101mm,height=44mm,xmin=-2.5,xmax=4.5,ymin=-1.5,ymax=3.5,xtick={-2,-1,0,1,2,3,4},ytick={-1,1,2,3}]
\addplot[otNavy,thick,mark=*]coordinates{(-2,1)(0,-1)(2,3)(4,1)};
\end{axis}\end{tikzpicture}
\end{center}
\begin{questions}
\item Résoudre $f(x)=1$ et $f(x)\geq1$.
\item Donner les variations et les extremums.
\end{questions}
\end{exercice}
\end{blocLivre}
\begin{blocLivre}{ex-14-7}
\begin{exercice}[2]{Des fréquences conditionnelles}
Dans un groupe de 100 personnes, 50 ont moins de 25 ans. Parmi ces 50, 20 utilisent le bus ; parmi les 50 autres, 30 l’utilisent.\begin{questions}
\item Calculer la fréquence globale d’utilisation du bus.
\item Calculer la fréquence d’utilisation parmi les moins de 25 ans.
\item Dans le modèle de tirage uniforme d’une personne, les deux évènements sont-ils indépendants ?
\end{questions}
\end{exercice}
\end{blocLivre}
\begin{blocLivre}{ex-14-8}
\begin{exercice}[2]{Réactiver la trigonométrie}
\begin{questions}
\item Donner les coordonnées du point-image de $5\pi/3$.
\item Comparer ses coordonnées à celles du point-image de $11\pi/3$.
\item Calculer un arc de rayon 3 et d’angle $2\pi/3$.
\end{questions}
\end{exercice}
\end{blocLivre}
\begin{blocLivre}{ex-14-9}
\begin{exercice}[2]{Relier dérivée et extremums}
$f(x)=x^3-6x^2+9x$ sur $[0;4]$.\begin{questions}
\item Factoriser $f\prime(x)$ et étudier son signe.
\item Donner le tableau de variations avec les valeurs aux bornes.
\item Donner toutes les positions des extremums globaux.
\end{questions}
\end{exercice}
\end{blocLivre}
\begin{blocLivre}{ex-14-10}
\begin{exercice}[2]{Un diagramme circulaire}
Lors d’un vote de 80 personnes, les options A, B et C recueillent respectivement 40, 24 et 16 voix.\begin{questions}
\item Calculer les fréquences et les angles des secteurs.
\item Quelle est la différence entre le nombre de voix pour A et pour B ?
\item Peut-on calculer une moyenne numérique des noms A, B et C ?
\end{questions}
\label{figlivre-r14-10-secteurs-ex}
\figureLivre{r14-10-secteurs}

\end{exercice}
\end{blocLivre}
\begin{blocLivre}{ex-14-11}
\begin{exercice}[2]{Une exponentielle dans une étude de fonction}
\begin{questions}
\item Résoudre $e^{2x}=e^{x+3}$.
\item Pour $g(x)=(2x-1)e^{-x}$ sur $[0;3]$, factoriser $g\prime(x)$.
\item Donner le maximum et le minimum de $g$ sur le segment.
\end{questions}
\end{exercice}
\end{blocLivre}
\clearpage
\rubrique{Problèmes}
\begin{blocLivre}{ex-14-12}
\begin{exercice}[3]{Synthèse : une boîte à dimensions liées}
Une boîte a une base carrée de côté $x$ et une hauteur $6-x$, en décimètres, avec $0<x<6$.\begin{questions}
\item Exprimer son volume $V(x)$ et calculer $V\prime(x)$.
\item Déterminer ses dimensions et son volume maximal.
\item Écrire la tangente à la courbe de $V$ en $x=2$.
\item Vérifier $V(x)-16=-(x-2)(x^2-4x-8)$, puis déterminer quand le volume est au moins 16 litres.
\end{questions}
\label{figlivre-r14-12-boite-ex}
\figureLivre{r14-12-boite}

\end{exercice}
\end{blocLivre}
\begin{blocLivre}{ex-14-13}
\begin{exercice}[3]{Synthèse : triangle, cercle et tangente}
On donne $A(0;0)$, $B(4;0)$ et $C(2;2)$.\begin{questions}
\item Montrer que le triangle est rectangle en $C$ et calculer l’angle $\widehat{BAC}$.
\item Déterminer le cercle circonscrit.
\item Écrire sa tangente en $C$ et projeter $M(5;3)$ sur cette tangente.
\item Calculer la longueur du petit arc $AC$ du cercle.
\end{questions}
\label{figlivre-r14-13-geometrie-ex}
\figureLivre{r14-13-geometrie}

\end{exercice}
\end{blocLivre}
\begin{blocLivre}{ex-14-14}
\begin{exercice}[3]{Synthèse : défauts et gain net}
Deux fournisseurs A et B livrent 60\,\% et 40\,\% des articles, avec des taux de défaut de 1\,\% et 4\,\%. Un article est vendu 20 euros, coûte 12 euros et entraîne 30 euros de dépense supplémentaire s’il est défectueux.\begin{questions}
\item Calculer la probabilité $P(D)$ d’un défaut.
\item Sachant un défaut, calculer la probabilité que l’article provienne de B.
\item Donner la loi du gain net G, puis son espérance et sa variance.
\item Pour quatre articles de défauts indépendants et identiquement distribués, calculer la probabilité d’au moins un défaut.
\end{questions}
\label{figlivre-r14-14-ex}
\figureLivre{r14-14}

\end{exercice}
\end{blocLivre}
\begin{blocLivre}{ex-14-15}
\begin{exercice}[3]{Synthèse : deux modèles d'évolution}
Un stock est modélisé par $u_0=1000$ et $u_{n+1}=0{,}9u_n+50$, au bout de $n$ semaines. On propose aussi le modèle continu $W(t)=500+500e^{-0{,}1t}$.\begin{questions}
\item Poser $v_n=u_n-500$ et en déduire une expression de $u_n$.
\item Justifier la décroissance de $u_n$ et montrer que $u_n>500$.
\item Par essais ou Python, trouver le premier rang tel que $u_n\leq600$.
\item Étudier les variations de W et trouver la première semaine entière où $W(n)\leq600$.
\item En utilisant $e^x\geq1+x$, avec égalité seulement en 0, comparer $u_n$ et $W(n)$ pour $n\geq1$. Expliquer la différence des seuils.
\end{questions}
\end{exercice}
\end{blocLivre}
